\documentclass{article}
\usepackage{/globals/de}
\usepackage{/globals/laws}
\usepackage{/globals/anki}
\ankiprefix{bundesverfassungsgericht}
\ankitopic{Bundesverfassungsgericht}
\ankishow
\begin{document}
\section{Das Bundesverfassungsgericht}
Das Bundesverfassungsgericht, das BVerfG, ist das oberste deutsche Gericht, mit dem Ziel die Verfassung zu schützen.

Es ist kein Fachgericht; es befasst sich nicht mit der gesamtheit des Rechts, sondern nur des Verfassungsrechts. Somit kann es \idr nicht für eine Steigerung einer Revision genutzt werden.
\ankinote{revision}{Revision}{Keine Steigerung einer Revision, keine Superrevisionsinstanz}

\arts{92}{94} \GG \gesetz{BVerfGG}

\subsection{Aufgaben}
Die zwei hauptsächlichen Aufgaben sind es Konflikte zwischen Verfassungsorganen zu schlichten und die Verfassungsrechte von Personen (Bürger, Unternehmen) zu schützen.
\ankinote{aufgaben}{Aufgaben}{Konflikte Verfassungsorgane, Verfassungsrechte Personen}

\subsection{Stellung}
Die Richter sind unabhängig (\arts{97} \absa{1} \GG).

Das BVerfG interpretiert und präzisiert Vorgaben aus dem Grundgesetz und bildet somit das Recht fort. Dies ist die \emph{richterliche Rechtsfortbildung}.
\ankinote{rechtsfortbildung}{Richterliche Rechtsfortbildung}{Interpretiert und präzisiert Grundgesetz}

Es hat das Recht Gesetze für ungültig zu erklären.
\ankinote{rechte}{Rechte}{Gesetze unnötig erklären, richterliche Rechtsfortbildung}

\subsection{Aufbau}
Das Gericht hat zwei Senate, je 8 Richter, davon je 3 ehemalige Bundesrichter. Sowohl der Bundestag als auch der Bundesrat wählen mit einer $2/3$ Mehrheit 8 je Richter.
\ankinote{aufbau}{Aufbau}{2 Senate, je 8 Richter, je 3 ehem. Bundesrichter}
\ankinote{wahl}{Wahl}{BT und BR wählen je 8 mit 2/3 Mehrheit}

\subsection{Verfahrensarten}
\begin{description}
 \item[Verfassungsbeschwerde] Alle Personen können Verfassungsbeschwerden einreichen wenn sie ihre Grundrechte als verletzt sehen.
 \item[Konkrete Normkontrolle] Es wird überprüft ob ein Bundesgesetz in einem konkreten Fall mit der Verfassung vereinbar ist.
 \item[Abstrakte Normkontrolle] Es wird geprüft ob ein Bundesgesetz ohne einem bestimmten Fall Verfassungsmäßig ist.
 \item[Organstreitverfahren] Konflikt zwischen zwei Bundesorganen über ihre Rechte und Pflichten.
 \item[Bund-Länder-Streitigkeiten] Konflikt zwischen dem Bund und einem Land über ihre Rechten und Pflichten.
\end{description}
\ankinote{va-arten}{Verfahrensarten Arten}{Verfassungsbeschwerde, Konkrete Normkontrolle, Abstrakte Normkontrolle, Organstreitverfahren, Bund-Länder-Streitigkeiten}
\ankinote{verfassungsbeschwerde}{Verfassungsbeschwerde}{Jede Person, Grundrechte verletzt}
\ankinote{konkrete-normkontrolle}{Konkrete Normkontrolle}{Bundesgesetz konkreter Fall Verfassungswidrig}
\ankinote{abstrakte-normkontrolle}{Abstrakte Normkontrolle}{Bundesgesetz nicht konkret Verfassungswidrig}
\ankinote{organstreitverfahren}{Organstreitverfahren}{Konflikte zwei Bundesorgane}
\ankinote{bund-länder-streitigkeiten}{Bund-Länder-Streitigkeiten}{Konflikt Bund, Land}

\end{document}